\chapter{Functions}

\section{Declaration}

Functions are declared with \lstinline{fn}:

\begin{lstlisting}[caption={Basic function}]
fn square(x) {
  return x^2
}

display square(5)    // 25
\end{lstlisting}

\subsection*{Implicit return}

The last expression in a function is returned automatically:

\begin{lstlisting}
fn square(x) {
  x^2    // implicit return
}
\end{lstlisting}

\subsection*{Multiple parameters}

\begin{lstlisting}
fn mean(a, b) {
  (a + b) / 2.0
}

display mean(10, 20)    // 15.0
\end{lstlisting}

\section{Functions as values}

Functions are first-class values — they can be assigned to variables,
passed as arguments, and returned by other functions:

\begin{lstlisting}[caption={Function as argument}]
fn apply(f, x) {
  f(x)
}

fn double(n) { n * 2 }

display apply(double, 7)    // 14
\end{lstlisting}

\section{Closures}

Closures are anonymous functions that capture the environment where they
were defined:

\begin{lstlisting}[caption={Basic closure}]
let multiplier = fn(factor) {
  fn(x) { x * factor }    // captures 'factor'
}

let triple = multiplier(3)
display triple(10)    // 30
\end{lstlisting}

\subsection*{Closures as arguments}

\begin{lstlisting}
fn apply_two(f, a, b) {
  f(a) + f(b)
}

let sum_of_squares = apply_two(fn(x) { x^2 }, 3, 4)
display sum_of_squares    // 9 + 16 = 25
\end{lstlisting}

\section{Recursion}

Functions can call themselves:

\begin{lstlisting}[caption={Recursive factorial}]
fn fact(n) {
  if n <= 1 { return 1 }
  n * fact(n - 1)
}

display fact(10)    // 3628800
\end{lstlisting}

\begin{nota}
  \hayashi{} does not optimise tail calls. For long sequences, prefer
  loops.
\end{nota}

\section{Optional parameters and defaults}

\hayashi{} has no native optional parameters, but the idiomatic pattern is
to use \lstinline{nil} as a sentinel and check explicitly:

\begin{lstlisting}
fn format_val(value, digits) {
  let d = if digits == nil { 2 } else { digits }
  // ...
}
\end{lstlisting}

\section{Function scope}

Functions declared with \lstinline{fn} at module level are available
throughout the file. Functions declared inside blocks are local to that
block:

\begin{lstlisting}
{
  fn local(x) { x + 1 }
  display local(5)    // 6
}
display local(5)    // error: undefined function 'local'
\end{lstlisting}

\section{\texttt{const} inside functions}

\lstinline{const} can be used inside functions to fix values that should
not change during execution:

\begin{lstlisting}
fn bmi(weight, height) {
  const FORMULA = "weight / height^2"
  weight / height^2
}
\end{lstlisting}
